% SPDX-License-Identifier: CC-BY-SA-4.0
% Copyright (c) 2026 Rafael Beraldo
% Parte 3: o espectro eletromagnético e a propagação
% Incluído por presentation.tex; compile a partir dele.

% =============================================================================
\section{O espectro eletromagnético}
% =============================================================================

\begin{frame}{Tudo é onda eletromagnética}
  \centering

    \foto[width=0.85\linewidth]{espectro-simulacao}{Rádio, micro-ondas, luz e raios X: a mesma onda, em frequências
    diferentes. Experimente:
    \linkexterno{\linkMiniPhysics}{Mini Physics, \emph{Electromagnetic spectrum explorer}}}
\end{frame}

\begin{frame}{Como a Anatel aloca o espectro de rádio brasileiro}
  \centering
  \figura{espectro-brasil}
\end{frame}

% Introdução à propagação (retomada depois com o perfil de uma travessia real).
\begin{frame}{Propagação: a frequência determina o alcance}
  \centering
  \figurapassos{propagacao}{3}

  \medskip
  \small
  \begin{tabular}{@{}rl@{}}
    \textcolor{cerjBlue}{\textbf{HF}} & reflete na ionosfera e pode ir longe, conforme a hora, a estação e o Sol \\
    \onslide<2->{\textcolor{cerjGreenDark}{\textbf{VHF/UHF}}} & \onslide<2->{não volta: só funciona com \alert{linha de visada}} \\
  \end{tabular}
\end{frame}

%%% Local Variables:
%%% TeX-master: "../presentation"
%%% fill-column: 100
%%% End:
